%!TeX root=../emmatop.tex
\chapter[Chapter \thechapter]{}
\lettrine[lines=4,lraise=0.3]{M}{rs} Weston's friends were all made happy by her safety; and if the satisfaction of her well-doing could be increased to Emma, it was by knowing her to be the mother of a little girl. She had been decided in wishing for a Miss Weston. She would not acknowledge that it was with any view of making a match for her, hereafter, with either of Isabella's sons; but she was convinced that a daughter would suit both father and mother best. It would be a great comfort to Mr Weston, as he grew older—and even Mr Weston might be growing older ten years hence—to have his fireside enlivened by the sports and the nonsense, the freaks and the fancies of a child never banished from home; and Mrs Weston—no one could doubt that a daughter would be most to her; and it would be quite a pity that any one who so well knew how to teach, should not have their powers in exercise again.

»She has had the advantage, you know, of practising on me,« she continued—»like La Baronne d'Almane on La Comtesse d'Ostalis, in Madame de Genlis' Adelaide and Theodore, and we shall now see her own little Adelaide educated on a more perfect plan.«

»That is,« replied Mr Knightley, »she will indulge her even more than she did you, and believe that she does not indulge her at all. It will be the only difference.«

»Poor child!« cried Emma; »at that rate, what will become of her?«

»Nothing very bad.—The fate of thousands. She will be disagreeable in infancy, and correct herself as she grows older. I am losing all my bitterness against spoilt children, my dearest Emma. I, who am owing all my happiness to you, would not it be horrible ingratitude in me to be severe on them?«

Emma laughed, and replied: »But I had the assistance of all your endeavours to counteract the indulgence of other people. I doubt whether my own sense would have corrected me without it.«

»Do you?—I have no doubt. Nature gave you understanding:—Miss Taylor gave you principles. You must have done well. My interference was quite as likely to do harm as good. It was very natural for you to say, what right has he to lecture me?—and I am afraid very natural for you to feel that it was done in a disagreeable manner. I do not believe I did you any good. The good was all to myself, by making you an object of the tenderest affection to me. I could not think about you so much without doating on you, faults and all; and by dint of fancying so many errors, have been in love with you ever since you were thirteen at least.«

»I am sure you were of use to me,« cried Emma. »I was very often influenced rightly by you—oftener than I would own at the time. I am very sure you did me good. And if poor little Anna Weston is to be spoiled, it will be the greatest humanity in you to do as much for her as you have done for me, except falling in love with her when she is thirteen.«

»How often, when you were a girl, have you said to me, with one of your saucy looks—»Mr Knightley, I am going to do so-and-so; papa says I may, or I have Miss Taylor's leave«—something which, you knew, I did not approve. In such cases my interference was giving you two bad feelings instead of one.«

»What an amiable creature I was!—No wonder you should hold my speeches in such affectionate remembrance.«

»»Mr Knightley.«—You always called me, »Mr Knightley;« and, from habit, it has not so very formal a sound.—And yet it is formal. I want you to call me something else, but I do not know what.«

»I remember once calling you »George,« in one of my amiable fits, about ten years ago. I did it because I thought it would offend you; but, as you made no objection, I never did it again.«

»And cannot you call me »George« now?«

»Impossible!—I never can call you any thing but »Mr Knightley.« I will not promise even to equal the elegant terseness of Mrs Elton, by calling you Mr K.—But I will promise,« she added presently, laughing and blushing—»I will promise to call you once by your Christian name. I do not say when, but perhaps you may guess where;—in the building in which N. takes M. for better, for worse.«

Emma grieved that she could not be more openly just to one important service which his better sense would have rendered her, to the advice which would have saved her from the worst of all her womanly follies—her wilful intimacy with Harriet Smith; but it was too tender a subject.—She could not enter on it.—Harriet was very seldom mentioned between them. This, on his side, might merely proceed from her not being thought of; but Emma was rather inclined to attribute it to delicacy, and a suspicion, from some appearances, that their friendship were declining. She was aware herself, that, parting under any other circumstances, they certainly should have corresponded more, and that her intelligence would not have rested, as it now almost wholly did, on Isabella's letters. He might observe that it was so. The pain of being obliged to practise concealment towards him, was very little inferior to the pain of having made Harriet unhappy.

Isabella sent quite as good an account of her visitor as could be expected; on her first arrival she had thought her out of spirits, which appeared perfectly natural, as there was a dentist to be consulted; but, since that business had been over, she did not appear to find Harriet different from what she had known her before.—Isabella, to be sure, was no very quick observer; yet if Harriet had not been equal to playing with the children, it would not have escaped her. Emma's comforts and hopes were most agreeably carried on, by Harriet's being to stay longer; her fortnight was likely to be a month at least. Mr and Mrs John Knightley were to come down in August, and she was invited to remain till they could bring her back.

»John does not even mention your friend,« said Mr Knightley. »Here is his answer, if you like to see it.«

It was the answer to the communication of his intended marriage. Emma accepted it with a very eager hand, with an impatience all alive to know what he would say about it, and not at all checked by hearing that her friend was unmentioned.

»John enters like a brother into my happiness,« continued Mr Knightley, »but he is no complimenter; and though I well know him to have, likewise, a most brotherly affection for you, he is so far from making flourishes, that any other young woman might think him rather cool in her praise. But I am not afraid of your seeing what he writes.«

»He writes like a sensible man,« replied Emma, when she had read the letter. »I honour his sincerity. It is very plain that he considers the good fortune of the engagement as all on my side, but that he is not without hope of my growing, in time, as worthy of your affection, as you think me already. Had he said any thing to bear a different construction, I should not have believed him.«

»My Emma, he means no such thing. He only means\longdash«

»He and I should differ very little in our estimation of the two,« interrupted she, with a sort of serious smile—»much less, perhaps, than he is aware of, if we could enter without ceremony or reserve on the subject.«

»Emma, my dear Emma\longdash«

»Oh!« she cried with more thorough gaiety, »if you fancy your brother does not do me justice, only wait till my dear father is in the secret, and hear his opinion. Depend upon it, he will be much farther from doing you justice. He will think all the happiness, all the advantage, on your side of the question; all the merit on mine. I wish I may not sink into »poor Emma« with him at once.—His tender compassion towards oppressed worth can go no farther.«

»Ah!« he cried, »I wish your father might be half as easily convinced as John will be, of our having every right that equal worth can give, to be happy together. I am amused by one part of John's letter—did you notice it?—where he says, that my information did not take him wholly by surprize, that he was rather in expectation of hearing something of the kind.«

»If I understand your brother, he only means so far as your having some thoughts of marrying. He had no idea of me. He seems perfectly unprepared for that.«

»Yes, yes—but I am amused that he should have seen so far into my feelings. What has he been judging by?—I am not conscious of any difference in my spirits or conversation that could prepare him at this time for my marrying any more than at another.—But it was so, I suppose. I dare say there was a difference when I was staying with them the other day. I believe I did not play with the children quite so much as usual. I remember one evening the poor boys saying, »Uncle seems always tired now.««

The time was coming when the news must spread farther, and other persons' reception of it tried. As soon as Mrs Weston was sufficiently recovered to admit Mr Woodhouse's visits, Emma having it in view that her gentle reasonings should be employed in the cause, resolved first to announce it at home, and then at Randalls.—But how to break it to her father at last!—She had bound herself to do it, in such an hour of Mr Knightley's absence, or when it came to the point her heart would have failed her, and she must have put it off; but Mr Knightley was to come at such a time, and follow up the beginning she was to make.—She was forced to speak, and to speak cheerfully too. She must not make it a more decided subject of misery to him, by a melancholy tone herself. She must not appear to think it a misfortune.—With all the spirits she could command, she prepared him first for something strange, and then, in a few words, said, that if his consent and approbation could be obtained—which, she trusted, would be attended with no difficulty, since it was a plan to promote the happiness of all—she and Mr Knightley meant to marry; by which means Hartfield would receive the constant addition of that person's company whom she knew he loved, next to his daughters and Mrs Weston, best in the world.

\begin{figure}[tbph]
\centering
\includegraphics[width=\linewidth]{53affectionately}
\caption{Emma hung about him affectionately}
\end{figure}

Poor man!—it was at first a considerable shock to him, and he tried earnestly to dissuade her from it. She was reminded, more than once, of having always said she would never marry, and assured that it would be a great deal better for her to remain single; and told of poor Isabella, and poor Miss Taylor.—But it would not do. Emma hung about him affectionately, and smiled, and said it must be so; and that he must not class her with Isabella and Mrs Weston, whose marriages taking them from Hartfield, had, indeed, made a melancholy change: but she was not going from Hartfield; she should be always there; she was introducing no change in their numbers or their comforts but for the better; and she was very sure that he would be a great deal the happier for having Mr Knightley always at hand, when he were once got used to the idea.—Did he not love Mr Knightley very much?—He would not deny that he did, she was sure.—Whom did he ever want to consult on business but Mr Knightley?—Who was so useful to him, who so ready to write his letters, who so glad to assist him?—Who so cheerful, so attentive, so attached to him?—Would not he like to have him always on the spot?—Yes. That was all very true. Mr Knightley could not be there too often; he should be glad to see him every day;—but they did see him every day as it was.—Why could not they go on as they had done?

Mr Woodhouse could not be soon reconciled; but the worst was overcome, the idea was given; time and continual repetition must do the rest.—To Emma's entreaties and assurances succeeded Mr Knightley's, whose fond praise of her gave the subject even a kind of welcome; and he was soon used to be talked to by each, on every fair occasion.—They had all the assistance which Isabella could give, by letters of the strongest approbation; and Mrs Weston was ready, on the first meeting, to consider the subject in the most serviceable light—first, as a settled, and, secondly, as a good one—well aware of the nearly equal importance of the two recommendations to Mr Woodhouse's mind.—It was agreed upon, as what was to be; and every body by whom he was used to be guided assuring him that it would be for his happiness; and having some feelings himself which almost admitted it, he began to think that some time or other—in another year or two, perhaps—it might not be so very bad if the marriage did take place.

Mrs Weston was acting no part, feigning no feelings in all that she said to him in favour of the event.—She had been extremely surprized, never more so, than when Emma first opened the affair to her; but she saw in it only increase of happiness to all, and had no scruple in urging him to the utmost.—She had such a regard for Mr Knightley, as to think he deserved even her dearest Emma; and it was in every respect so proper, suitable, and unexceptionable a connexion, and in one respect, one point of the highest importance, so peculiarly eligible, so singularly fortunate, that now it seemed as if Emma could not safely have attached herself to any other creature, and that she had herself been the stupidest of beings in not having thought of it, and wished it long ago.—How very few of those men in a rank of life to address Emma would have renounced their own home for Hartfield! And who but Mr Knightley could know and bear with Mr Woodhouse, so as to make such an arrangement desirable!—The difficulty of disposing of poor Mr Woodhouse had been always felt in her husband's plans and her own, for a marriage between Frank and Emma. How to settle the claims of Enscombe and Hartfield had been a continual impediment—less acknowledged by Mr Weston than by herself—but even he had never been able to finish the subject better than by saying—»Those matters will take care of themselves; the young people will find a way.« But here there was nothing to be shifted off in a wild speculation on the future. It was all right, all open, all equal. No sacrifice on any side worth the name. It was a union of the highest promise of felicity in itself, and without one real, rational difficulty to oppose or delay it.

Mrs Weston, with her baby on her knee, indulging in such reflections as these, was one of the happiest women in the world. If any thing could increase her delight, it was perceiving that the baby would soon have outgrown its first set of caps.

The news was universally a surprize wherever it spread; and Mr Weston had his five minutes share of it; but five minutes were enough to familiarise the idea to his quickness of mind.—He saw the advantages of the match, and rejoiced in them with all the constancy of his wife; but the wonder of it was very soon nothing; and by the end of an hour he was not far from believing that he had always foreseen it.

»It is to be a secret, I conclude,« said he. »These matters are always a secret, till it is found out that every body knows them. Only let me be told when I may speak out.—I wonder whether Jane has any suspicion.«

\begin{figure}[tbph]
\centering
\includegraphics[width=.9\linewidth]{53passed}
\caption{It passed to Mrs Cole, Mrs Perry, and Mrs Elton}
\end{figure}

He went to Highbury the next morning, and satisfied himself on that point. He told her the news. Was not she like a daughter, his eldest daughter?—he must tell her; and Miss Bates being present, it passed, of course, to Mrs Cole, Mrs Perry, and Mrs Elton, immediately afterwards. It was no more than the principals were prepared for; they had calculated from the time of its being known at Randalls, how soon it would be over Highbury; and were thinking of themselves, as the evening wonder in many a family circle, with great sagacity.

In general, it was a very well approved match. Some might think him, and others might think her, the most in luck. One set might recommend their all removing to Donwell, and leaving Hartfield for the John Knightleys; and another might predict disagreements among their servants; but yet, upon the whole, there was no serious objection raised, except in one habitation, the Vicarage.—There, the surprize was not softened by any satisfaction. Mr Elton cared little about it, compared with his wife; he only hoped »the young lady's pride would now be contented;« and supposed »she had always meant to catch Knightley if she could;« and, on the point of living at Hartfield, could daringly exclaim, »Rather he than I!«—But Mrs Elton was very much discomposed indeed.—»Poor Knightley! poor fellow!—sad business for him.«—She was extremely concerned; for, though very eccentric, he had a thousand good qualities.—How could he be so taken in?—Did not think him at all in love—not in the least.—Poor Knightley!—There would be an end of all pleasant intercourse with him.—How happy he had been to come and dine with them whenever they asked him! But that would be all over now.—Poor fellow!—No more exploring parties to Donwell made for her. Oh! no; there would be a Mrs Knightley to throw cold water on every thing.—Extremely disagreeable! But she was not at all sorry that she had abused the housekeeper the other day.—Shocking plan, living together. It would never do. She knew a family near Maple Grove who had tried it, and been obliged to separate before the end of the first quarter.